\documentclass[11pt]{article}
\usepackage[margin=1in]{geometry}
\usepackage{parskip}
\begin{document}

\begin{center}
{\LARGE Effect of Daily Stair Climbing on Resting Heart Rate:\\ A 6-Week Randomized Controlled Trial in Office Workers}\\[6pt]
{\large A. Researcher, B. Author, and C. Collaborator}\\[4pt]
{\small Department of Exercise Science, Northfield University}
\end{center}

\section*{Abstract}

Sedentary office work is associated with poor cardiovascular fitness. We conducted a randomized controlled trial of 64 office workers (32 intervention, 32 control) to test whether a simple stair-climbing intervention reduces resting heart rate, a marker of cardiovascular fitness. The intervention group was instructed to climb at least five flights of stairs each workday for six weeks; the control group maintained their usual routine. The primary outcome was resting heart rate measured at week 6. Resting heart rate decreased more in the intervention group than in the control group (mean difference -3.1 beats per minute, 95\% CI -5.4 to -0.8, p = 0.009), with a between-group effect size of d = 0.38. Adherence to the stair-climbing instruction was 82\%. The trial was pre-registered and funded by the university.

\section*{Methods}

We recruited 64 healthy adults from two office buildings and randomized them 1:1 to a six-week stair-climbing intervention or a usual-activity control condition. The intervention group received a daily prompt to climb at least five flights of stairs. Resting heart rate was measured at baseline and week 6 using an automated monitor after five minutes of seated rest. The analysis used an analysis of covariance adjusting for baseline heart rate, age, and sex. Adherence was tracked by self-report.

\section*{Results}

Baseline resting heart rate was 72.4 bpm in the intervention group and 71.9 bpm in the control group. At week 6, the intervention group had a mean resting heart rate of 67.1 bpm and the control group 70.2 bpm (adjusted mean difference -3.1 bpm, 95\% CI -5.4 to -0.8, p = 0.009). The between-group effect size was d = 0.38, a small-to-moderate effect. Adherence to the stair-climbing instruction was 82\% overall.

\section*{Discussion}

Six weeks of daily stair climbing produced a statistically significant but modest reduction in resting heart rate relative to usual activity. The effect size was small-to-moderate. Strengths include the randomized design, pre-registration, and adjustment for baseline covariates. Limitations include a short follow-up period, self-reported adherence, and a sample restricted to healthy office workers, which limits generalizability. Longer trials with objective activity monitoring are needed to determine whether the effect is durable and clinically meaningful.

\end{document}
